\documentclass[11pt]{article}
\usepackage{amsmath, amssymb, amsfonts}
\usepackage{geometry}
\geometry{letterpaper, margin=1in}
\usepackage{booktabs}

\newcommand{\NameText}{\texttt{SO(n=4, q=2)}}
\newcommand{\MatrixSize}{4}
\newcommand{\Rank}{2}

\title{Pinned Group Computation Report: \NameText{}}
\author{Pinning-Sympy Automated Output}
\date{\today}

\begin{document}

\maketitle
\tableofcontents

\newpage

\section{Basic attributes}

\begin{tabular}{ll}
    \toprule
    Group & \NameText \\
    \midrule
    Matrix size & $\MatrixSize \times \MatrixSize$ \\
    Rank & $\Rank$ \\
    \bottomrule
\end{tabular}

\subsection{Form matrix}
\[
\left[\begin{matrix}0 & 0 & 1 & 0\\0 & 0 & 0 & 1\\1 & 0 & 0 & 0\\0 & 1 & 0 & 0\end{matrix}\right]
\]

\subsection{Generic Lie algebra element}
\[
\left[\begin{matrix}x_{0 0} & x_{0 1} & 0 & - x_{1 2}\\x_{1 0} & x_{1 1} & x_{1 2} & 0\\0 & - x_{3 0} & - x_{0 0} & - x_{1 0}\\x_{3 0} & 0 & - x_{0 1} & - x_{1 1}\end{matrix}\right]
\]

\subsection{Generic torus element}
\[
\left[\begin{matrix}t_{0} & 0 & 0 & 0\\0 & t_{1} & 0 & 0\\0 & 0 & \frac{1}{t_{0}} & 0\\0 & 0 & 0 & \frac{1}{t_{1}}\end{matrix}\right]
\]

\subsection{Trivial characters}
\[
\left[\begin{matrix}1 & 0 & 1 & 0\\0 & 1 & 0 & 1\\1 & 1 & 1 & 1\end{matrix}\right]
\]

\newpage
\section{Root system}

\subsection*{Summary Properties}
\begin{tabular}{ll}
    \toprule
    \textbf{Property} & \textbf{Value} \\
    \midrule
    Dynkin type & \texttt{A1\_x\_A1} \\
    Reduced & True \\
    Simply laced & True \\
    Number of roots & 4 \\
    \bottomrule
\end{tabular}

\subsection*{Decomposition Components}
This is a reducible root system with the following components:

\noindent \begin{tabular}{cc}
    \toprule
    \textbf{Component Type} & \textbf{Component Rank} \\
    \midrule
    \texttt{A1} & 1 \\
    \texttt{A1} & 1 \\
    \bottomrule
\end{tabular}

\subsection*{Root to Coroot Mapping}
\begin{center}
\begin{tabular}{cc}
    \toprule
    \textbf{Root ($\alpha$)} & \textbf{Coroot ($\alpha^{\vee}$)} \\
    \midrule
    $\left( -1, \  -1, \  0, \  0\right)$ & $\left( -1, \  -1, \  1, \  1\right)$ \\
    $\left( -1, \  1, \  0, \  0\right)$ & $\left( -1, \  1, \  1, \  -1\right)$ \\
    $\left( 1, \  -1, \  0, \  0\right)$ & $\left( 1, \  -1, \  -1, \  1\right)$ \\
    $\left( 1, \  1, \  0, \  0\right)$ & $\left( 1, \  1, \  -1, \  -1\right)$ \\
    \bottomrule
\end{tabular}
\end{center}


\end{document}